\section{大域モジュライへの適用条件}\label{sec:component}
\begin{theorem}[相対モジュライ比較を明示した成分判定]\label{thm:conditional-component}
次の条件を仮定する。
\begin{enumerate}
\item 問題のH3対象を含む局所有限型のrigidified代数的モジュライstack $\mathfrak M$ があり、
上記の平坦な反標準環の族を受け入れる。
\item 一般点でH3の無限小変形から構造対象付き圏の無限小変形への写像が単射である。
\item 第\ref{sec:integrate}節の八次元一般表示族とその七次元族への特殊化が、
この相対モジュライ問題の族として実現される。
\end{enumerate}
このとき七次元族の一般点を通る既約成分は唯一で、次元は8である。
その点で $\mathfrak M$ は滑らかで、完成局所底は八次元形式円板となる。
\end{theorem}
\begin{proof}
条件(2)と定理\ref{thm:HH8}から接空間次元は8以下である。
条件(3)の一般表示族は、境界の平行移動を動かし、同型による冗長性を除いて八次元である。
その閉包が七次元族の一般点を含むので、局所次元は8以上となる。
一般慣性は自明であるため\'{e}tale局所環を取れ、接空間次元とKrull次元はともに8。
従って局所環は正則であり、唯一の既約成分を持つ。
\end{proof}

第\ref{sec:artin-reconstruction}節で条件(2)を適合したenhanced Abel変形の範囲では証明する。
他の条件を未検証のまま外して「全成分を完全に決定した」とはしない。
特に正則な八次元形式近傍を得ても、全既約成分の境界の記述はそれだけでは得られない。
完全決定には、一般族の像の閉包に含まれるH3対象を漏れなく認識し、
滑らかな楕円境界が消える点、中心が特異になる点、非PI領域、
別の表示を持つ点について全矢と慣性を求める必要がある。

有力な候補は、Etingof--Ginzburg型の一般中心三次商を含む八次元成分である。
Rainsの $E_6$ 型幾何学と比較する際には、
固定した楕円境界のfiltered deformationの六次元と、境界曲線・平行移動の二次元を区別する。
期待される重み付き射影表示やWeyl商を、そのまま全H3の同値と断定しない。
指定potential chartでの全Artin包絡認識と八変数の同時中心切断は後節で解決する。
全intrinsic形式変形との比較、H3族の相対実現、閉包の境界の認識は保持すべき課題である。
他方、指標多様体・特異点・実幾何との具体的接続は、それらの完成を待たずに調べられる。
